%% Manuscript template: numbers in double-brace tokens are filled from results/ by fill_manuscript.py.
%% Edit this file, then run: python fill_manuscript.py  -> ../rGO_CuO_JVP_Piezoresistive paper/JVP_ML_manuscript.tex
\documentclass[11pt]{article}
\usepackage[a4paper,margin=2.3cm]{geometry}
\usepackage[T1]{fontenc}
\usepackage{newtxtext,newtxmath}
\usepackage{amsmath}
\usepackage{graphicx}
\usepackage{booktabs,tabularx,multirow}
\usepackage{xcolor}
\usepackage[super,sort&compress,comma]{natbib}
\usepackage[hidelinks]{hyperref}
\usepackage{lineno}
\usepackage{setspace}
\usepackage{caption}
\captionsetup{font=small,labelfont=bf,labelsep=period}
\graphicspath{ {figures_jvp/} }
\newcommand{\authnote}[1]{\textcolor{red}{[\textbf{AUTHORS:} #1]}}
\renewcommand{\figurename}{Fig.}
\linenumbers
\onehalfspacing

\title{\textbf{A wearable rGO/CuO piezoresistive jugular venous pulse sensor with physiology-informed machine learning for screening heart disease, benchmarked against the electrocardiogram}}
\author{Arijit Pattra$^{1}$, Ayan Chatterjee$^{1}$, Sayan Dey$^{1,*}$\\[4pt]
\small $^{1}$School of Electrical and Computer Sciences, Indian Institute of Technology Bhubaneswar,\\ \small Argul, Khordha, Odisha 752050, India\\
\small $^{*}$Corresponding author: sayandey@iitbbs.ac.in}
\date{}

\begin{document}
\maketitle

\begin{abstract}
\noindent The jugular venous pulse (JVP) is the body's natural manometer of right-heart haemodynamics, yet it is still assessed by subjective bedside inspection or by invasive catheterisation. Here we present a skin-mounted reduced graphene oxide/copper oxide (rGO/CuO) piezoresistive sensor that records the JVP waveform over the external jugular vein, and we test whether machine learning on the sensor signal alone, as it would be used in deployment, can identify heart disease; electrocardiogram (ECG) leads V1 and V2 served only as a benchmark. In {{N}} participants ({{NPOS}} with physician-verified heart disease; {{NCYC}} sensor windows), all sensor analyses were start-agnostic, because the position of each recorded window within the cardiac cycle was undocumented. Under subject-level, repeated, nested cross-validation, the best sensor model, a shift-invariant one-dimensional convolutional neural network, reached an AUROC of {{ACI:jvp_cnn}}, and sparse logistic regression {{ACI:jvp_l1}} (permutation $P={{L1PERM_P}}$), compared with {{ACI:ecg_best}} for the best ECG V1/V2 comparator model (CNN vs ECG, DeLong {{DLP:cnn_ecglr}}). The hypothesised ratio of \textit{v}-wave to \textit{c}-wave amplitude was only weakly discriminative (AUROC {{VCAUC}}), whereas an exploratory shift-invariant measure of waveform asymmetry separated the groups (AUROC {{AUC_REL3}}). An alignment test showed that models that are not start-agnostic reach near-perfect apparent accuracy (AUROC {{AL:stored|raw}}) that falls to chance after random circular shifts of the windows ({{AL:shift|raw}}). Class-imbalance corrections mainly altered calibration and operating points. With {{NPOS}} diseased participants and a model selected among eleven, these results indicate that the sensor alone carries a moderate signal, comparable to a two-lead ECG, and motivate a prospective study with continuous, synchronised recordings and about {{SS:auc0.85|n_pos}} diseased participants.
\end{abstract}

\noindent\textbf{Keywords:} jugular venous pulse; wearable sensor; reduced graphene oxide; piezoresistive; machine learning; right ventricular hypertrophy; heart block; electrocardiogram

%% ======================================================================
\section*{Introduction}

Cardiovascular disease (CVD) remains the leading cause of death worldwide. Prevalent CVD nearly doubled to 523 million cases between 1990 and 2019 \cite{roth2020global}, CVD deaths reached 19.8 million in 2022 \cite{mensah2023global}, and heart failure (HF) alone affects more than 64 million people \cite{savarese2023global}. The burden is disproportionately borne by low- and middle-income countries. In India, CVD caused 28.1\% of all deaths in 2016, more than half of them before the age of 70 \cite{indiastatelevel2018changing,prabhakaran2016cardiovascular}, and patients hospitalised with HF are younger and die in hospital more often than those in Western registries \cite{harikrishnan2015clinical}. Right-heart pressure and volume overload are common contributors: pulmonary hypertension affects about 1\% of the world's population, roughly 80\% of whom live in developing countries \cite{hoeper2016global,humbert2022esc}, and rheumatic heart disease, a major cause of right-sided valve disease, remains concentrated in South Asia \cite{watkins2017global}. Conduction disease is also consequential; first-degree atrioventricular (AV) block roughly triples the long-term risk of pacemaker implantation \cite{cheng2009longterm,kerola2019risk}, and high-grade block requires pacing \cite{kusumoto2019acc,glikson2021esc}. Detecting these conditions currently depends on ECG and echocardiography interpreted by specialists, yet India has only about five to six active qualified doctors per 10,000 people, unevenly distributed \cite{karan2021size}. Low-cost, objective and wearable tools for first-line screening are therefore a priority.

The jugular venous pulse offers such a window. Because no valve separates the internal jugular vein from the right atrium, the neck-vein pulsation follows right-atrial pressure and right-heart haemodynamics \cite{applefeld1990jugular,chuachiaco2013jugular,constant2000using,devine2007jugular}. The normal waveform comprises the \textit{a} wave (atrial contraction, following the ECG P wave), the \textit{c} wave and \textit{x} descent (tricuspid closure, ventricular systole and atrial relaxation), the \textit{v} wave (passive atrial filling against a closed tricuspid valve, peaking near the end of the T wave), and the \textit{y} descent (tricuspid opening) \cite{applefeld1990jugular,chuachiaco2013jugular,ranganathan2023jugular}. Each component is time-locked to an electrical event, so the JVP encodes both mechanics and conduction. Prominent \textit{a} waves indicate resistance to right-ventricular (RV) filling, as in RV hypertrophy (RVH) and pulmonary hypertension; cannon \textit{a} waves mark AV dissociation such as complete heart block; and large \textit{v} or \textit{cv} waves indicate tricuspid regurgitation and RV volume overload \cite{applefeld1990jugular,devine2007jugular,mcgee2018inspection,morin2017extreme,tung2016cannon,mansoor2015kussmaul}. The JVP also carries prognosis: in 2,479 patients from SOLVD, an elevated JVP independently predicted HF hospitalisation and pump-failure death \cite{drazner2001prognostic}, and jugular estimates of right-atrial pressure track left-sided filling pressures \cite{drazner2008value,butman1993bedside}. It therefore remains a cornerstone of congestion assessment in guidelines \cite{mcdonagh2021esc,gheorghiade2010assessing,thibodeau2018role}.

In practice, however, the bedside JVP is unreliable. Its examination is subjective and posture-dependent, inter-observer agreement is poor, and clinicians underestimate catheter-measured venous pressure \cite{cook1990clinical,cook1996does,mcgee1998physical}. The conventional sternal-angle reference varies widely between patients \cite{seth2002far}; physical signs together were only 58\% sensitive for elevated filling pressures \cite{stevenson1989limited}; and residents detected raised right-atrial pressure with 14\% sensitivity, compared with 82\% for hand-held ultrasound \cite{brennan2007comparison}. The jugular vein cannot be assessed at all in up to 30\% of patients \cite{vinayak2006usefulness,sankoff2008noninvasive}. Ultrasound of the internal jugular vein improves estimation \cite{simon2010detection,pellicori2014revisiting,pellicori2015prognostic,ammirati2024estimation,rudski2010guidelines}, but requires equipment and a trained operator and provides a snapshot rather than the continuous \textit{a--c--x--v--y} time course. Electrocardiographic screening has complementary limits: ECG criteria for RVH \cite{sokolow1949ventricular,hancock2009aha} are specific but insensitive, with 18--44\% sensitivity in autopsy series \cite{murphy1984reevaluation,lehtonen1988electrocardiographic} and a positive predictive value of at most 12\% against cardiac MRI in MESA-RV \cite{whitman2014validity}. Implantable pressure sensors show that continuous haemodynamic information changes outcomes, reducing HF hospitalisation in CHAMPION \cite{abraham2011wireless} and, in its pre-COVID analysis, GUIDE-HF \cite{lindenfeld2021haemodynamic,zile2008transition,stevenson2010chronic}, but they are invasive and costly.

\paragraph{Related work: non-invasive and wearable JVP sensing.}
Several sensing routes to the JVP have been explored (Table~\ref{tab:prior}). Non-contact optical methods were pioneered by Amelard et al., who showed that photoplethysmographic imaging of the neck reveals a JVP waveform anti-correlated with the carotid pulse \cite{amelard2017noncontact}, that optical jugular venous attenuation tracks invasive central venous pressure \cite{amelard2021optical}, and that it can quantify jugular compliance \cite{amelard2022assessing}. Skin-displacement imaging \cite{lampotang2018noncontact}, camera vibrocardiography \cite{moco2017camera,moco2018importance}, Eulerian video magnification against right-heart catheterisation \cite{abnousi2019novel}, remote video examination \cite{kelly2020feasibility} and smartphone video \cite{farahani2025noncontact} followed, as reviewed in ref.~\cite{arrow2023capturing}; these approaches require a still subject, controlled lighting and a camera in front of the neck \cite{saiko2024observation}. Contact optical sensing was demonstrated with anterior-neck reflectance photoplethysmography referenced to ECG and ultrasound \cite{garcialopez2020extracting,garcialopez2020characterization}, and machine learning was recently used to separate venous from arterial beats in wearable neck PPG \cite{cao2026unmixing}. Ultrasonic approaches range from a conformal ultrasonic patch that captured jugular waveforms \cite{wang2018monitoring}, to wearable cardiac imaging and Doppler \cite{hu2023wearable,sempionatto2021epidermal,wang2021flexible,kenny2021novel}, to single-element A-mode jugular tracking in 65 volunteers \cite{george2022high,george2025single,sahani2016carotid}. Radio-frequency near-field coherent sensing \cite{hui2018monitoring} produced a collar-mounted JVP sensor with ECG- and PCG-consistent wave timings \cite{conroy2023jugular}, and 60-GHz radar has been shown to be feasible \cite{das2025feasibility}. Mechanical skin-contact sensors are closest to the present device: cervical strain-gauge plethysmography \cite{proto2021plethysmography,menegatti2022effect}, a gyroscope patch \cite{karhinoja2023method} and a contact piezoelectric sensor that agreed with ultrasound ($r=0.90$) in 20 volunteers \cite{george2024jugular}. Flexible e-skin devices, including a microhair capacitive sensor \cite{pang2015highly}, a carbon-nanotube/PDMS piezoresistive film \cite{tai2015flexible} and a piezoelectret neck array \cite{han2022mapping}, have recorded the jugular pulse only as demonstrations.

\paragraph{Related work: flexible piezoresistive materials.}
Flexible pressure sensing grew from microstructured dielectrics and transistors \cite{schwartz2013flexible,wang2014silk}, interlocking nanostructures \cite{pang2012flexible}, nanowire and piezoelectric films \cite{gong2014wearable,dagdeviren2014conformable,park2017self} and bioresorbable arterial sensors \cite{boutry2019biodegradable}, as reviewed in refs.~\cite{chortos2016pursuing,ray2019bio,meng2022wearable}. Graphene and rGO piezoresistors combine high gauge factors with low-cost solution processing: graphene--polyurethane sponges \cite{yao2013flexible}, laser-scribed and spinosum-mimetic graphene \cite{tian2015graphene,pang2018epidermis}, rGO paper and rGO-coated textiles \cite{tao2017graphene,li2019wearable}, MXene and MXene/rGO aerogels \cite{ma2017highly,ma2018synergistical} and rGO-wrapped copper nanowire networks \cite{zhu2022epidermis}. These devices are almost always demonstrated on the radial or carotid arterial pulse, whose amplitude is an order of magnitude larger than that of the venous pulse. To our knowledge, an rGO/CuO hybrid has not been reported as a piezoresistive venous-pulse sensor.

\paragraph{Related work: machine learning on cardiac waveforms.}
Deep learning has transformed ECG interpretation, from cardiologist-level rhythm classification with a 1D convolutional neural network (CNN) trained on 91,232 recordings \cite{hannun2019cardiologist} and 12-lead classification on two million tracings \cite{ribeiro2020automatic}, to the detection of conditions invisible to human readers, including LV systolic dysfunction \cite{attia2019screening,yao2021artificial}, occult atrial fibrillation \cite{attia2019artificial}, LV hypertrophy \cite{kwon2020comparing}, pulmonary hypertension \cite{kwon2020artificial,aras2023electrocardiogram}, RV dysfunction \cite{vaid2022using} and incident complete heart block \cite{sau2025artificial} (reviewed in refs.~\cite{siontis2021artificial,somani2021deep,topol2019high}; benchmarks in refs.~\cite{wagner2020ptbxl,strodthoff2021deep}). Machine learning on mechanical and haemodynamic waveforms is less mature: smartwatch PPG detects atrial fibrillation \cite{tison2018passive,perez2019large} and even diabetes \cite{avram2020digital}; arterial pulse-wave analysis underpins cuffless blood pressure estimation \cite{mukkamala2015toward,elhajj2020review,kim2020detection}; and seismocardiography with machine learning assesses HF status and intracardiac pressures \cite{inan2015ballistocardiography,inan2018novel,shandhi2022estimation,klein2025noninvasive,etemadi2016wearable,stehlik2020continuous,lee2020mechano,ha2019chest,kaisti2019clinical}. These successes rely on thousands to millions of labelled recordings. For the JVP, machine learning has served only signal processing---source separation in neck PPG \cite{cao2026unmixing} and segmentation of jugular ultrasound \cite{zamboni2026artificial}---rather than diagnosis.

\paragraph{Gap and contributions.}
Thus, although the JVP has been recorded with optical, ultrasonic, radio-frequency and mechanical sensors, almost all studies report waveform morphology or timing agreement in small groups of healthy volunteers. We found no wearable, skin-conformal JVP sensor that has been combined with machine-learning classification of heart disease and benchmarked against the ECG. Here we address this gap with three contributions. (i) We use an rGO/CuO piezoresistive sensor placed over the external jugular vein to record JVP windows in {{N}} participants alongside reference ECG leads V1 and V2, and analyse them with methods that make no assumption about where each window starts. (ii) We translate clinical hypotheses---that RV disease raises the \textit{v} wave relative to the \textit{c} wave, and that conduction delay flattens the \textit{a}$\to$\textit{c} transition---into automatically measured waveform features, and test them statistically. (iii) Treating the sensor as a standalone device, with the ECG used only as a comparator, we benchmark interpretable models, random-kernel and convolutional neural networks, and a physiology-hybrid network against ECG-based models under subject-level repeated nested cross-validation, with bootstrap confidence intervals, permutation tests and DeLong comparisons, following TRIPOD+AI recommendations \cite{collins2024tripod}, and we analyse explicitly how the 1:8 class imbalance affects discrimination, operating points and calibration. Because the cohort contains only {{NPOS}} diseased participants, we frame the study as a feasibility and proof-of-concept analysis and quantify its uncertainty explicitly \cite{varoquaux2018cross,riley2020calculating}.

\begin{table}[t]
\centering\small
\caption{\textbf{Prior non-invasive JVP sensing approaches.} Subject numbers are as reported in each abstract. None of the prior approaches performed disease classification.}
\label{tab:prior}
\begin{tabularx}{\textwidth}{>{\raggedright}p{3.9cm}>{\raggedright}p{1.5cm}>{\raggedright}p{1.2cm}>{\raggedright\arraybackslash}X p{0.9cm}}
\toprule
Modality & Contact & Wearable & Validation reference (cohort) & Ref. \\
\midrule
Camera PPG imaging (neck) & Non-contact & No & Carotid PPG, ultrasound (healthy) & \citenum{amelard2017noncontact} \\
Optical jugular venous attenuation & Non-contact & No & Invasive CVP (healthy) & \citenum{amelard2021optical} \\
Camera skin displacement & Non-contact & No & Laser displacement, ECG (19 healthy) & \citenum{lampotang2018noncontact} \\
Eulerian video magnification & Non-contact & No & Right-heart catheter (50 patients; clinician read) & \citenum{abnousi2019novel} \\
Anterior-neck reflectance PPG & Contact & Yes & ECG, B-mode ultrasound (20 healthy) & \citenum{garcialopez2020extracting} \\
Dual-channel neck PPG + $k$NN & Contact & Yes & Venous vs arterial beats (15 healthy) & \citenum{cao2026unmixing} \\
Conformal ultrasonic patch & Contact & Yes & Reference devices (healthy) & \citenum{wang2018monitoring} \\
Single-element ultrasound & Contact & Portable & Ultrasound imaging (65 volunteers) & \citenum{george2025single} \\
Near-field RF collar & Proximal & Yes & ECG, PCG timing (21) & \citenum{conroy2023jugular} \\
60-GHz FMCW radar & Non-contact & No & Neck PPG; one catheter trace & \citenum{das2025feasibility} \\
Cervical strain-gauge plethysmography & Contact & Yes & ECG, ultrasound (20 healthy) & \citenum{proto2021plethysmography},\citenum{menegatti2022effect} \\
Gyroscope patch & Contact & Yes & ECG, SCG (20 healthy) & \citenum{karhinoja2023method} \\
Contact piezoelectric sensor & Contact & Semi & Ultrasound, $r=0.90$ (20 healthy) & \citenum{george2024jugular} \\
Microhair capacitive / CNT piezoresistive film & Contact & Yes & Waveform demonstration & \citenum{pang2015highly},\citenum{tai2015flexible} \\
\midrule
\textbf{This work: rGO/CuO piezoresistive + ML} & Contact & Yes & ECG V1/V2 ({{N}} participants, {{NPOS}} diseased); \textbf{disease classification} & -- \\
\bottomrule
\end{tabularx}
\end{table}

%% ======================================================================
\section*{Results}

\subsection*{Sensor, cohort and data curation}
\authnote{Insert device description and characterisation here (rGO/CuO synthesis; morphology by SEM, XRD, Raman and XPS; electrode geometry, encapsulation, sensitivity/gauge factor, linear range, response/recovery time, hysteresis, cyclic stability, temperature and humidity drift, force calibration). The current draft does not contain these data; the figures in this folder from the As(III) FET paper are not applicable to this device.}

Recordings from {{N}} participants (subject folders HS1--HS54; HS10 and HS14 absent) were available, labelled at the subject level as diseased ($n={{NPOS}}$) or not diseased ($n={{NNEG}}$) on the basis of physician diagnosis, which was independently verified \authnote{specify the diagnostic work-up (e.g. 12-lead ECG, echocardiography), the number and specialty of the physicians, whether they had access to the sensor recordings, the specific diagnoses (RVH, heart block, other) of the {{NPOS}} diseased participants, inclusion/exclusion criteria, age, sex and BMI, and ethics approval}. Each participant contributed two to four short sensor windows (sensor current and calibrated force; 0.6--1.1\,s, about one cardiac period each) and one reference beat each from ECG leads V1 and V2 (Fig.~\ref{fig:pipeline}). After removing exact duplicates, {{NCYC}} unique windows remained (median {{SPC}} samples per window, range {{SPC_RANGE}}). Because the position of each window within the cardiac cycle is not documented, we made no assumption about where a window starts and did not use window duration; all sensor analyses were designed to be invariant to the window's starting point (Methods). The force channel was a per-recording linear transform of the current (median $|r|={{RIF}}$), so we analysed force and verified that raw current gave the same results.

A systematic audit of the workbooks (Supplementary Table~S1) found issues that are typical of manually assembled datasets and that we handled explicitly: the ECG beats of four participants (HS42, HS43, HS48, HS52) were identical to those of other participants, and the headers of three workbooks (HS31, HS43, HS48) named a different participant, consistent with copied templates; ECG time stamps in three participants were stored in the wrong unit (milliseconds, a factor-of-10 error, or an implausible scale) and were repaired; single-lead workbook headers sometimes named the wrong lead (lead identity was taken from the file name); four participants had a cycle duplicated under two cycle numbers; and one JVP cycle differed between workbooks of the same participant. All analyses used the curated data, and a sensitivity analysis excluded the five participants with copied workbooks.

\begin{figure}[t]
\centering
\includegraphics[width=\textwidth]{fig1_pipeline.pdf}
\caption{\textbf{Study pipeline.} JVP windows recorded by the rGO/CuO piezoresistive sensor over the external jugular vein (EJV) and reference ECG leads V1/V2 are processed into start-agnostic (shift-invariant) features and trough-anchored waveforms. Sensor-only models (the intended deployed system) are compared with ECG-only comparator models under subject-level repeated stratified cross-validation with nested hyper-parameter selection; no model combines the two modalities.}
\label{fig:pipeline}
\end{figure}

\subsection*{Sensor waveform morphology}
Because the position of each sensor window within the cardiac cycle is unknown, every window was rotated so that its deepest trough, taken as the \textit{x} descent, lies at the origin; the prominent peaks then follow in the physiological order \textit{x}$\to$\textit{v}$\to$\textit{y}$\to$\textit{a}$\to$\textit{c} (Fig.~\ref{fig:waves}a,b). The class-mean waveforms overlapped widely (Fig.~\ref{fig:waves}c): non-diseased windows declined smoothly from the \textit{v} wave towards the next \textit{x} trough, whereas diseased windows tended to show a deeper trough after the \textit{v} wave (\textit{y} descent) followed by a distinct late \textit{c} wave. The class-mean ECG in V1 of diseased participants showed a deeper S wave (Fig.~\ref{fig:waves}d).

\begin{figure}[t]
\centering
\includegraphics[width=\textwidth]{fig2_waveforms.pdf}
\caption{\textbf{JVP windows recorded by the rGO/CuO sensor.} \textbf{a,b}, Trough-anchored windows of a non-diseased (HS13) and a diseased (HS5) participant (64-point monotone cubic interpolation); the first window is annotated with the automatically detected \textit{x}, \textit{v}, \textit{y}, \textit{a} and \textit{c} fiducials. \textbf{c}, Class-mean trough-anchored window ($\pm$s.d.). \textbf{d}, Class-mean reference ECG beat in lead V1 ($\pm$s.d.).}
\label{fig:waves}
\end{figure}

\subsection*{Testing the physiological hypotheses}
The \textit{v}/\textit{c} amplitude ratio was higher in diseased participants (median {{VC_MED_DIS}} vs {{VC_MED_CTL}}), but the difference was not significant (Mann--Whitney $P={{VCP}}$; AUROC {{VCAUC}}, 95\% bootstrap CI {{VCCI}}; Fig.~\ref{fig:features}a). The hypothesised rule \textit{v}/\textit{c}$>$1 identified {{RULE_TP}} of {{NPOS}} diseased participants (sensitivity {{RULE_SENS}}) with a specificity of {{RULE_SPEC}} ({{RULE_TN}} of {{NNEG}}). The \textit{a}$\to$\textit{c} slope did not separate the groups (AUROC {{AUC_SLOPE}}, $P={{P_SLOPE}}$). In an exploratory screen of all 33 shift-invariant sensor features, the strongest separation came from the sine of the third-harmonic relative phase, a descriptor of how the large and small waves are arranged within the period (AUROC {{AUC_REL3}}, $P={{P_REL3}}$; significant after Bonferroni correction for 33 features), followed by the depth of the \textit{y} descent (AUROC {{AUC_YD}}, $P={{P_YD}}$). On the ECG, diseased participants had a deeper S wave in V1 (AUROC {{AUC_V1S}}), consistent with the hypothesis of a sharp S wave, whereas the standard ECG voltage criteria for RVH (R in V1 $\geq$0.7\,mV or R/S in V1 $\geq$1 \cite{hancock2009aha}) flagged only {{RVH_TP}} of {{NPOS}} diseased participants (sensitivity {{RVH_SENS}}, specificity {{RVH_SPEC}}), in line with their known insensitivity \cite{murphy1984reevaluation,whitman2014validity}.

\begin{figure}[t]
\centering
\includegraphics[width=\textwidth]{fig3_features.pdf}
\caption{\textbf{Testing the physiological hypotheses.} Subject-level values (median over windows) for non-diseased ($n={{NNEG}}$) and diseased ($n={{NPOS}}$) participants. \textbf{a}, \textit{v}/\textit{c} amplitude ratio; dashed line, the hypothesised threshold of 1. \textbf{b}, \textit{a}$\to$\textit{c} slope. \textbf{c}, Sine of the third-harmonic relative phase (waveform asymmetry; exploratory). \textbf{d}, Depth of the \textit{y} descent. \textbf{e}, ECG S-wave amplitude in V1. Horizontal bars, medians. Titles give the AUROC (orientation-free) and two-sided Mann--Whitney $P$, unadjusted.}
\label{fig:features}
\end{figure}

\subsection*{Benchmarking sensor-based models against ECG}
Table~\ref{tab:bench} and Fig.~\ref{fig:bench} summarise all models under subject-level repeated stratified three-fold cross-validation (each test fold contains two diseased participants), with hyper-parameters and decision thresholds selected inside each training fold. Among the eleven sensor models, the start-agnostic 1D-CNN performed best (AUROC {{ACI:jvp_cnn}}; AUPRC {{PR:jvp_cnn}} against a chance level of {{PREV}}), followed by sparse logistic regression ({{ACI:jvp_l1}}), the physiology-hybrid CNN ({{ACI:jvp_hcnn}}), circular ROCKET ({{ACI:jvp_rocket}}) and logistic regression on the Fourier descriptors ({{ACI:jvp_fourier}}). Models that used all 33 features without selection performed poorly (logistic regression {{ACI:jvp_lr}}; permutation $P={{PERM:jvp_lr}}$), consistent with overfitting when the features outnumber the diseased participants more than five-fold, and the \textit{v}/\textit{c} rule reached {{ACI:rule}}. Because the 1D-CNN was identified as the best of eleven sensor models after the comparison, its estimate is optimistic.

The best ECG model, logistic regression on V1/V2 morphology features, reached {{ACI:ecg_best}}, and random forest {{ACI:ecg_rf}}; an ECG CNN with a single beat per participant reached {{ACI:ecg_cnn}}. The sensor 1D-CNN and the best ECG model did not differ significantly ({{DLTXT:cnn_ecglr}}, DeLong), and their screening decisions agreed only {{KAPPACNNTXT}} (Cohen's $\kappa={{KAPPACNN}}$). In permutation tests that re-ran the complete cross-validation on permuted labels, sparse logistic regression on sensor features ($P={{L1PERM_P}}$, {{L1PERM_N}} permutations) and the ECG model ($P={{PERM:ecg_lr}}$) performed above chance. For the 1D-CNN, the test had to use a reduced configuration (one cross-validation pass and one seed; Methods), under which its AUROC was lower ({{CNNPERM_OBS}}) and the result was not significant ($P={{CNNPERM_P}}$, {{CNNPERM_N}} permutations); its above-chance performance is therefore supported only indirectly, by the significant result of the similarly performing sparse model. The \textit{v}/\textit{c}-only model did not exceed chance ($P={{PERM:jvp_lrvc}}$).

\begin{table}[t]
\centering\footnotesize
\caption{\textbf{Subject-level cross-validated performance} ({{N}} participants, {{NPOS}} diseased; prevalence {{PREV}}). AUROC and AUPRC are means over cross-validation repeats, with 95\% stratified subject-bootstrap CIs. Sensitivity, specificity, balanced accuracy and the Matthews correlation coefficient (MCC) use thresholds chosen inside each training fold (0.5 for the CNNs and ROCKET); the AUPRC of a random classifier equals the prevalence. LR, logistic regression; L1-LR, sparse logistic regression; SVM, support vector machine; RF, random forest; GBM, gradient boosting; CNN, convolutional neural network.}
\label{tab:bench}
\resizebox{\textwidth}{!}{%
\begin{tabular}{llccccccc}
\toprule
Modality & Model & AUROC (95\% CI) & AUPRC (95\% CI) & Sens. & Spec. & Bal. acc. & MCC & Brier \\
\midrule
{{BENCHTABLE}}
\bottomrule
\end{tabular}}
\end{table}

\begin{figure}[t]
\centering
\includegraphics[width=\textwidth]{fig4_benchmark.pdf}
\caption{\textbf{Sensor versus ECG benchmark.} \textbf{a}, Cross-validated AUROC with 95\% bootstrap CIs for all sensor (violet; the deployed modality) and ECG V1/V2 (green; comparator only) models; dotted line, chance. \textbf{b}, ROC and \textbf{c}, precision--recall curves of representative models using scores averaged over cross-validation repeats (legend values: AUROC and average precision). The grey line in \textbf{c} marks the prevalence ({{PREV}}).}
\label{fig:bench}
\end{figure}

\subsection*{What the models learn}
Fitted on all participants, sparse logistic regression retained a single feature, the sine of the third-harmonic relative phase, with lower values favouring disease (Fig.~\ref{fig:interp}a); all other coefficients were shrunk to zero. The gradient$\times$input saliency of the 1D-CNN was diffuse (Fig.~\ref{fig:interp}b): for diseased windows it was positive over the \textit{v}-wave segment and peaked in the late segment before the return to \textit{x}, where non-diseased windows had negative saliency, indicating that the network contrasted the shape of the late part of the beat, in keeping with the class-mean waveforms. Saliency maps from a model trained on six diseased participants should be interpreted qualitatively.

\begin{figure}[t]
\centering
\includegraphics[width=\textwidth]{fig5_interpretability_robustness.pdf}
\caption{\textbf{Interpretability and robustness.} \textbf{a}, Standardised coefficients of sparse logistic regression fitted on all participants (ten largest in magnitude; all but one are zero); blue, lower values favour disease. \textbf{b}, Gradient$\times$input saliency of the sensor 1D-CNN averaged over windows of each class, in trough-anchored coordinates; vertical lines mark the median fiducial phases. \textbf{c}, Permutation tests; vertical lines, observed AUROC (*for the CNN, under the reduced configuration used for its null distribution); shading spans the null median to its 95th percentile. \textbf{d}, AUROC of models trained on clean data and tested on signals with added Gaussian noise (SNR 30 to 0\,dB) or decimated to one-half or one-third of the samples (mean and 95\% interval over 10 repeats).}
\label{fig:interp}
\end{figure}

\subsection*{Robustness and sensitivity analyses}
When models trained on clean signals were tested on noisy signals (Fig.~\ref{fig:interp}d), the 1D-CNN kept an AUROC of {{DEG:SNR 10 dB|CNN}} at 10\,dB SNR, falling to {{DEG:SNR 5 dB|CNN}} at 5\,dB and {{DEG:SNR 0 dB|CNN}} at 0\,dB; sparse logistic regression gave {{DEG:SNR 10 dB|LR}}, {{DEG:SNR 5 dB|LR}} and {{DEG:SNR 0 dB|LR}}. Reducing the sampling density to one-half or one-third changed the CNN little ({{DEG:1/2 samples|CNN}} and {{DEG:1/3 samples|CNN}}). Excluding the five participants with copied workbooks gave {{SENS:QC-flagged workbooks excluded|JVP 1D-CNN}} for the 1D-CNN and {{SENS:QC-flagged workbooks excluded|ECG LR V1/V2}} for ECG logistic regression, and using raw current instead of calibrated force gave {{SENS:raw current instead of force|JVP 1D-CNN}}. With a single window per participant, the 1D-CNN fell to {{SENS:single cycle per subject|JVP 1D-CNN}}, indicating that averaging over several beats matters (Supplementary Fig.~S1).

\subsection*{Accounting for class imbalance}
With {{NPOS}} diseased among {{N}} participants (prevalence {{PREVPCT}}\%), a classifier that calls everyone non-diseased is {{ACCALLNEG}}\% accurate, so we report only imbalance-robust metrics---AUROC, AUPRC (chance level {{PREV}}), Matthews correlation coefficient (MCC) and balanced accuracy \cite{saito2015precision,chicco2020advantages}---and tested whether correcting the imbalance changes the conclusions (Supplementary Table~S2 and Fig.~S3). All corrections were applied inside the training folds only, including the inner tuning folds, so that synthetic or re-weighted samples never reached test participants. Discrimination changed comparatively little with the correction strategy: sensor logistic regression on all features reached an AUROC of {{IMB:JVP features|LR|none|auroc}} without correction, {{IMB:JVP features|LR|class weight|auroc}} with class weighting, {{IMB:JVP features|LR|SMOTE|auroc}} with SMOTE \cite{chawla2002smote} and {{IMB:JVP features|LR|undersampling ensemble|auroc}} with an undersampling ensemble \cite{liu2009exploratory} (random forest {{IMBMIN:JVP features|RF|auroc}}--{{IMBMAX:JVP features|RF|auroc}}), so resampling, if anything, reduced the discrimination of the sensor feature model. ECG logistic regression gained slightly from correction ({{IMBMIN:ECG features|LR|auroc}}--{{IMBMAX:ECG features|LR|auroc}}) and remained above the sensor feature models in every configuration.

The imbalance instead determined the operating point and the calibration. Without correction, the default threshold of 0.5 detected few diseased participants (sensitivity {{IMB:JVP features|LR|none|sens@0.5}} for sensor and {{IMB:ECG features|LR|none|sens@0.5}} for ECG logistic regression), even when the ranking was good. Class weighting, SMOTE and undersampling raised the sensor model's sensitivity at 0.5 to {{IMBS05MIN}}--{{IMBS05MAX}}, and choosing the threshold inside the training folds achieved a similar operating point without any correction (sensitivity {{IMB:JVP features|LR|none|sens@youden}}, specificity {{IMB:JVP features|LR|none|spec@youden}}). The price of the corrections was miscalibration: they raised the mean predicted risk {{IMBCITLMIN}}--{{IMBCITLMAX}} above the observed prevalence, compared with {{IMB:JVP features|LR|none|citl}} without correction, in line with the known harms of imbalance corrections for risk prediction \cite{vandengoorbergh2022harm}. The sensor 1D-CNN behaved in the same way: its AUROC was {{IMBMIN:JVP waveform|1D-CNN|auroc}}--{{IMBMAX:JVP waveform|1D-CNN|auroc}} with unweighted, positive-weighted, focal-loss \cite{lin2017focal} or oversampled training, whereas its sensitivity at 0.5 was {{IMB:JVP waveform|1D-CNN|none|sens@0.5}} without correction and {{IMB:JVP waveform|1D-CNN|weighted|sens@0.5}} with positive weighting.

At clinically relevant operating points, the sensor 1D-CNN had a sensitivity of {{OP:JVP 1D-CNN|sens@spec90}} at 90\% specificity and {{OP:JVP 1D-CNN|sens@spec95}} at 95\% specificity, compared with {{OP:ECG LR|sens@spec90}} and {{OP:ECG LR|sens@spec95}} for ECG logistic regression. Because predictive values depend on prevalence, we projected them to screening settings (Fig.~S3e): at a prevalence of 5\%, the PPV would be {{PV:JVP 1D-CNN|0.05|ppv}} for the sensor 1D-CNN and {{PV:ECG LR|0.05|ppv}} for ECG logistic regression, with NPVs of {{PV:JVP 1D-CNN|0.05|npv}} and {{PV:ECG LR|0.05|npv}}. Finally, the small number of events limits precision. For a true AUROC of 0.85, close to the values observed here, about {{SS:auc0.85|n_pos}} diseased participants ({{SS:auc0.85|n_total}} in total at the same prevalence) would be needed to estimate it within $\pm$0.05 \cite{hanley1982meaning} (about $\pm${{SSHW}} for an AUROC of 0.90 with the present {{NPOS}}), and at least {{SSSENS}} diseased participants to estimate a sensitivity of 0.9 within $\pm$0.1.

\subsection*{Dependence on window alignment}
The stored windows are not randomly positioned within the cardiac cycle: their deepest trough lay at a median of {{TROUGHPOS}} of the window (interquartile range {{TROUGHIQR}}), earlier in diseased ({{TROUGHDIS}}) than in non-diseased participants ({{TROUGHCTL}}). Any model that is not start-agnostic can therefore exploit where in the window the waves happen to fall. Logistic regression on the raw window reached an AUROC of {{AL:stored|raw}} on the stored windows but {{AL:shift|raw}}, close to chance, after each window was randomly circularly shifted, whereas logistic regression on shift-invariant Fourier descriptors was unaffected ({{AL:stored|fourier}} and {{AL:shift|fourier}}; Fig.~\ref{fig:agree}a). Positional information of this kind cannot be interpreted physiologically unless the windows are extracted from the sensor signal by a documented, pre-specified rule, so all primary analyses were start-agnostic. Coupling between start-agnostic sensor features and ECG features was weak (Fig.~\ref{fig:agree}b): only two of 63 pairs reached nominal significance, the third-harmonic relative phase with the V1 S-wave amplitude ($\rho={{RHO:relphase3_sin|V1_S_mV}}$; lower phase values accompanying deeper S waves) and the \textit{v}/\textit{c} ratio with the V2 R-wave amplitude ($\rho={{RHO:vc_ratio|V2_R_mV}}$), and neither survives correction for multiple comparisons.

\begin{figure}[t]
\centering
\includegraphics[width=\textwidth]{fig6_ecg_agreement.pdf}
\caption{\textbf{Alignment dependence and coupling with the ECG.} \textbf{a}, Cross-validated AUROC (mean and 95\% interval over 20 repeats) of logistic regression on the raw 64-point window and on shift-invariant Fourier descriptors, for the windows as stored and after a random circular shift of every window; dotted line, chance. \textbf{b}, Spearman correlations between start-agnostic sensor features and ECG features; values shown where $P<0.05$ (unadjusted). Participants with copied ECG beats are excluded.}
\label{fig:agree}
\end{figure}

\subsection*{Interpreting the ECG benchmark}
ECG interpretation is the clinical standard for diagnosing conduction disease and a routine screen for RVH, and published deep-learning ECG models detect pulmonary hypertension and RV dysfunction with AUROCs of 0.84--0.90 \cite{kwon2020artificial,aras2023electrocardiogram,vaid2022using}. The ECG AUROC of {{A:ecg_best}} observed here is in line with those figures despite an impoverished input (Supplementary Fig.~S2). First, the ECG files are not recorded signals: a six-Gaussian P--Q--R--S--S$'$--T model reproduced them with a median $R^2$ of {{EI:r2}} ({{EI:r2n}} of {{EI:nleads}} leads $>$0.999), and a median {{EI:flat}}\% of each beat was exactly flat. Second, the reconstructed beats carry no conduction information: the PR interval was below the physiological minimum of 120\,ms in {{EI:pr120}} of {{N}} participants and exceeded the 200\,ms threshold for first-degree AV block in {{EI:pr200}} of {{N}}, so AV conduction delay could not be detected from these beats by any model. Third, the ECG models were trained on 52 single beats of two leads, whereas state-of-the-art ECG networks are trained on $10^5$--$10^6$ full 12-lead recordings \cite{hannun2019cardiologist,attia2019screening,ribeiro2020automatic}, and the classical RVH voltage criteria rely on leads (V5, V6, limb leads for axis) that were not available \cite{hancock2009aha}. A crude start-agnostic \textit{v}/\textit{c} ratio from the raw sensor samples agreed with the diagnosis in {{EI:agree}} of {{N}} participants (AUROC {{EI:auc}}). Both modalities therefore carry a moderate signal in these data, and neither comparison establishes how the sensor would perform against a clinical 12-lead ECG.

%% ======================================================================
\section*{Discussion}
We show that a skin-mounted rGO/CuO piezoresistive sensor records jugular venous pulse waveforms from which start-agnostic models distinguish participants with physician-verified heart disease from those without with moderate accuracy: the best sensor model, a 1D-CNN, reached an AUROC of {{A:jvp_cnn}} and a sparse logistic regression {{A:jvp_l1}} (permutation $P={{L1PERM_P}}$), not significantly different from the best ECG V1/V2 model ({{A:ecg_best}}). To our knowledge, this is the first wearable JVP sensor to be combined with machine-learning classification of heart disease and benchmarked against the ECG; earlier wearable and non-contact JVP studies reported waveform or timing agreement in healthy volunteers \cite{amelard2017noncontact,garcialopez2020extracting,conroy2023jugular,karhinoja2023method,george2024jugular,george2025single} or used machine learning for signal separation \cite{cao2026unmixing}.

The most important methodological finding is that the apparent diagnostic accuracy of JVP analyses can depend heavily on how the recorded windows are positioned within the cardiac cycle. A model that reads the raw window reached near-perfect discrimination on the stored windows, but this advantage disappeared when the windows were randomly rotated, whereas genuinely shift-invariant representations were unaffected; any analysis that locates waves at fixed positions within the window is exposed to the same effect. Positional information is not necessarily spurious, because the timing of the venous waves relative to atrial and ventricular events is physiologically meaningful \cite{applefeld1990jugular,chuachiaco2013jugular}; however, it can be used only if the windows are extracted from the sensor signal by a documented rule, or synchronised to a simultaneously recorded ECG, and it must be validated as such. We therefore recommend that wearable JVP studies report their segmentation procedure and include an alignment test of the kind reported here.

The physiological picture is more nuanced than the original hypotheses. Once alignment information was removed, the \textit{v}/\textit{c} ratio alone was only weakly discriminative, and the \textit{a}$\to$\textit{c} slope did not separate the groups, although diseased participants tended to have higher \textit{v}/\textit{c} ratios and deeper \textit{y} descents, as expected with RV volume overload and tricuspid regurgitation \cite{applefeld1990jugular,devine2007jugular,ranganathan2023jugular}. The strongest single descriptor was a measure of waveform asymmetry, suggesting that the arrangement of the venous waves within the beat, rather than any single amplitude ratio, carries the discriminative information; because it emerged from an exploratory screen, it requires confirmation. The best models were those that limit effective complexity, the 1D-CNN with shared convolutional weights and sparse logistic regression, whereas unregularised use of all 33 features overfitted; this agrees with evidence that model complexity must be matched to the number of events in small clinical datasets \cite{christodoulou2019systematic,grinsztajn2022why,riley2020calculating}.

The class-imbalance analysis carries a practical lesson for deployment: re-weighting or resampling was not needed to rank participants correctly, but the decision threshold must be chosen deliberately for the intended prevalence, and models trained with imbalance corrections must be recalibrated before their outputs are read as risks \cite{vandengoorbergh2022harm}. The ECG comparison should be interpreted with care. The ECG inputs were single parametric reconstructions of two leads without conduction information, the V1/V2 models were trained on 52 beats, and standard voltage criteria for RVH are known to be insensitive \cite{murphy1984reevaluation,lehtonen1988electrocardiographic,whitman2014validity}. The weak coupling between sensor and ECG features suggests that the two modalities capture partly different information; because the sensor is intended to be used on its own, the ECG served here only as a benchmark of what conventional electrical screening achieves in the same participants.

\paragraph{Limitations.} This study has important limitations, and its results should be regarded as proof of concept. First, the cohort is small and highly imbalanced: with six diseased participants, the confidence intervals are wide and the estimates may not generalise \cite{varoquaux2018cross,vabalas2019machine,riley2020calculating}. Second, the analysis was revised after inspection of the results: the start-agnostic design, the trough-anchored fiducials and two of the models (sparse logistic regression and the Fourier-descriptor model) were introduced in a re-analysis, and the 1D-CNN was selected as the primary sensor model because it performed best among eleven, so its AUROC is optimistic and requires confirmation in independent data. Third, the diseased group pools different conditions (RVH and heart block) \authnote{confirm diagnoses}, so we cannot report condition-specific performance, and a pilot cohort of clear-cut patients and volunteers may separate more easily than unselected patients (spectrum bias) \cite{wolff2019probast,bossuyt2015stard}. Fourth, the sensor windows were short (one beat), sampled at only about 5--25\,Hz and of undocumented position within the cardiac cycle, and the ECG beats were reconstructed rather than recorded synchronously, which precluded timing analyses; ECG intervals measured automatically on these reconstructions included some implausible values. Fifth, several workbooks contained copy errors; excluding them did not change the conclusions. Finally, all data come from one site, without external validation. A prospective, multi-centre study with continuous, synchronised sensor and 12-lead ECG acquisition at $\geq$100\,Hz, automated sensor-only segmentation specified in advance, echocardiographic reference standards, condition-specific labels and blinded assessment, reported according to TRIPOD+AI and STARD-AI \cite{collins2024tripod,sounderajah2025stard,bossuyt2015stard,wolff2019probast}, is the necessary next step \cite{kelly2019key,liu2020reporting}.

In conclusion, an rGO/CuO piezoresistive neck sensor converts the jugular venous pulse, a subjective bedside sign, into a quantitative waveform that, analysed without any assumption about window alignment, separated participants with heart disease from those without with moderate accuracy, comparable to that of a two-lead ECG, in this pilot cohort; the estimates are imprecise and require confirmation. Low-cost wearable JVP sensing with start-agnostic machine learning merits prospective validation for heart-disease screening where specialists and echocardiography are scarce.

%% ======================================================================
\section*{Methods}

\subsection*{Sensor fabrication and characterisation}
\authnote{Describe the rGO/CuO synthesis (e.g. GO by modified Hummers' method; CuO formation route; reduction conditions; composition ratio), film deposition, substrate, electrodes and encapsulation; materials characterisation; electromechanical characterisation (current--force calibration, sensitivity, detection limit, linearity, response/recovery, hysteresis, cycling stability, temperature and humidity effects); read-out electronics (bias voltage, source-measure unit, sampling rate). The text in achemso-demo.tex describes an rGO/GO FET arsenic sensor and must not be reused.}

\subsection*{Participants and data acquisition}
\authnote{Recruitment setting and dates; inclusion and exclusion criteria; reference-standard diagnosis (physician-verified; specify modality, number of physicians and whether they were blinded to the sensor data); ethics committee approval number and informed consent; participant position (e.g. supine at 30--45$^\circ$), sensor placement over the right external jugular vein, fixation, breathing instructions, recording duration; ECG device, leads and sampling rate; how the ECG beats in the ``ECG\_Coordinates'' files were obtained (digitised from printed ECG? reconstructed?).} \authnote{state how the sensor windows were extracted from the continuous recording (automatically or manually; with or without reference to the ECG).} For each participant the dataset contained two to four sensor windows (time, current in A, calibrated force in N), one ECG beat for V1 and one for V2 (time, amplitude in mV), and a subject-level label.

\subsection*{Intended use}
The system under evaluation is the sensor alone: at deployment, a participant would be screened from sensor readings without an ECG. The ECG leads V1 and V2 were therefore used only to build comparator models that indicate what conventional electrical screening achieves in the same participants, and never as inputs to, or for the alignment of, the sensor models. The primary sensor model, trained on all participants with operating thresholds taken from its cross-validated scores, is provided as a sensor-only research tool with the code (Code availability).

\subsection*{Data curation}
Workbooks were parsed with column positions located from the header text, because positions varied between files. Sensor windows repeated across the V1, V2 and V1\_V2 workbooks of a participant were de-duplicated by content hash. ECG lead identity was taken from the file name. ECG time stamps were repaired when the beat duration was implausible (stored in milliseconds; a factor-of-10 error; or a near-zero duration, in which case the duration of the other lead was used). Byte-identical ECG beats across participants and workbook headers naming another participant were flagged (Supplementary Table~S1). No participant was excluded from the primary analysis.

\subsection*{Signal processing and features}
Only the sensor signal was used to derive sensor features; ECG timing was never used to locate JVP waves. Each sensor window was interpolated onto 64 points with monotone piecewise cubic (PCHIP) interpolation \cite{fritsch1980monotone}, which avoids overshoot between the sparse samples, and min--max normalised to remove inter-individual differences in contact pressure and sensor gain. Because the window's starting point within the cardiac cycle is unknown, each window was treated as one period of a periodic signal and every feature was made invariant to a circular shift. \textit{Trough-anchored fiducials:} the window was rotated so that its deepest trough, taken as the \textit{x} descent, was at the origin; prominent peaks (prominence $\geq$0.08 of the range, detected circularly) then follow in physiological order \textit{x}$\to$\textit{v}$\to$(\textit{y})$\to$\textit{a}$\to$\textit{c}, so \textit{v} was the first prominent peak after \textit{x}, \textit{c} the last before the return to \textit{x}, \textit{a} the peak preceding \textit{c} (when three or more peaks were present) and \textit{y} the minimum between \textit{v} and the next peak; with a single dominant hump, \textit{v} and \textit{c} were the maxima of the first and second halves after \textit{x}. From these we computed wave heights above \textit{x}, the ratios \textit{v}/\textit{c}, \textit{v}/\textit{a} and \textit{c}/\textit{a} (each regularised by adding 0.05 to numerator and denominator), the \textit{y} descent, phase distances (\textit{x}$\to$\textit{v}, \textit{v}$\to$\textit{c}, \textit{c}$\to$\textit{x}), the \textit{a}$\to$\textit{c} slope, \textit{v} upstroke and \textit{x} descent slopes, and the number of prominent peaks. \textit{Fourier descriptors:} normalised magnitudes of harmonics 1--6 and the relative phases $\phi_k-k\phi_1$ ($k=2,3,4$; as cosine and sine), which are invariant to circular shifts and encode waveform shape and asymmetry. \textit{Distribution and texture:} skewness, kurtosis, fraction of the window above half-range and circular roughness. \textit{Beat-to-beat measures:} s.d. of \textit{v}/\textit{c} and mean maximum circular cross-correlation between windows. The 33 features were aggregated per participant by the median over windows. ECG features for V1 and V2 were baseline-corrected R and S amplitudes (within 60\,ms of the steepest QRS slope), R/S ratio, R$-|$S$|$, QRS duration (onset at the PR-segment minimum preceding the first QRS deflection; offset at the return to within 15\% of the QRS amplitude), P and T amplitudes, PR interval (P onset at 10\% of the P amplitude), QT interval and beat length, together with the RVH voltage criteria \cite{hancock2009aha}. Feature extraction used no labels, so it could not leak label information across folds.

\subsection*{Models}
\textit{Rule.} The hypothesised rule (diseased if \textit{v}/\textit{c}$>$1) was applied without training. \textit{Tabular models.} Logistic regression (L2, class-balanced), support vector machine (RBF kernel, ranked by margin), random forest (500 trees) \cite{breiman2001random} and gradient boosting \cite{friedman2001greedy}, each with median imputation and standardisation, were tuned by an inner stratified three-fold grid search on AUROC \cite{cawley2010overfitting,vabalas2019machine} using scikit-learn \cite{pedregosa2011scikit}. The decision threshold was chosen by the Youden index on inner out-of-fold scores. Sparse (L1) logistic regression, which selects features inside each training fold, and L2 logistic regression restricted to the 12 Fourier descriptors were added during the start-agnostic re-analysis because the full feature set has far more features than diseased participants. \textit{ROCKET.} One thousand random convolutional kernels \cite{dempster2020rocket} were applied circularly to each trough-anchored window, so that the proportion-of-positive-values and maximum features are exactly invariant to circular shifts, followed by strongly regularised logistic regression; participant scores averaged window scores. \textit{1D-CNN.} A three-block convolutional network with circular padding (16--32--32 channels; kernels 7, 5, 3; batch normalisation; global average pooling; dropout 0.3; about 9,000 parameters) took the normalised window and its circular derivative as two channels \cite{kiranyaz2021onedimensional,ismailfawaz2019deep}. Windows were randomly circularly shifted during training, and predictions were averaged over eight evenly spaced circular shifts at test time, making the network start-agnostic. It was trained for 120 epochs with AdamW \cite{loshchilov2019decoupled} (learning rate $3\times10^{-3}$, weight decay $10^{-3}$, cosine schedule), positive-class-weighted binary cross-entropy, and on-the-fly augmentation by smooth time warping, amplitude scaling, jitter and baseline wander \cite{um2017data,iwana2021empirical}. Three seeds were ensembled, and participant probability was the mean over windows. \textit{Physiology-hybrid CNN.} The CNN embedding was concatenated with the two standardised features that encode the a-priori clinical hypotheses (\textit{v}/\textit{c} ratio and \textit{a}$\to$\textit{c} slope) before the output layer. \textit{ECG models} used the same tabular learners on V1/V2 features and a CNN on the resampled V1/V2 beat (two channels, 200 samples). No model combined sensor and ECG inputs, because the intended system uses the sensor alone; the ECG models serve only as a benchmark. Models were implemented in Python with scikit-learn and PyTorch \cite{paszke2019pytorch}.

\subsection*{Class imbalance}
Primary models used class-balanced weights (tabular learners) or positive-class-weighted cross-entropy (CNNs), with thresholds chosen inside the training folds. In a dedicated analysis we compared, for logistic regression and random forests on sensor and ECG features, no correction, class weighting, SMOTE \cite{chawla2002smote} (synthetic positives interpolated between each positive and one of its three nearest positive neighbours until the classes were balanced) and an undersampling ensemble (ten models, each trained on all positives and an equal number of randomly drawn negatives) \cite{liu2009exploratory}; for the sensor 1D-CNN we compared unweighted and positive-weighted cross-entropy, class-balanced focal loss ($\gamma=2$) \cite{lin2017focal} and oversampling of positive windows to parity. Resampling was implemented inside the estimators so that it was applied within both outer and inner folds and never to test data \cite{he2009learning}. We report AUROC, AUPRC, MCC \cite{chicco2020advantages}, sensitivity and specificity at a threshold of 0.5 and at the in-fold Youden threshold, Brier score and calibration-in-the-large (mean predicted risk minus prevalence) \cite{vancalster2019calibration}. Sensitivity at fixed specificity and specificity at 100\% sensitivity were computed on the out-of-fold scores of each repeat and averaged. Positive and negative predictive values at other prevalences were derived from the cross-validated sensitivity and specificity with Bayes' theorem, and sample sizes from the Hanley--McNeil standard error of the AUROC \cite{hanley1982meaning,riley2020calculating}.

\subsection*{Evaluation and statistics}
All splits were made at the participant level to prevent leakage between cycles of the same person \cite{saeb2017need,kapoor2023leakage}. We used repeated stratified three-fold cross-validation (30 repeats for fast models and 10 for CNN, ROCKET and tree models), with all tuning and threshold selection nested within the training folds. For each repeat, out-of-fold predictions for all participants were pooled to compute AUROC, average precision (AUPRC; appropriate under class imbalance \cite{saito2015precision}), sensitivity, specificity, balanced accuracy and Brier score. Point estimates are means over repeats, and 95\% CIs are from 2,000 stratified participant-level bootstrap resamples of that mean \cite{efron1979bootstrap,efron1993introduction}. Paired AUROC comparisons used the DeLong test on repeat-averaged scores \cite{delong1988comparing}. Statistical significance against chance was assessed with permutation tests that re-ran the complete cross-validation on {{NPERM}} label permutations \cite{ojala2010permutation}; for the 1D-CNN, whose full configuration is too costly to re-run hundreds of times, the observed statistic and the null distribution ({{CNNPERM_N}} permutations) were both computed with one cross-validation pass and one seed. Univariate group differences used two-sided Mann--Whitney tests and are reported without multiplicity adjustment as exploratory. To quantify how much an analysis that is not start-agnostic could depend on where the windows happen to start, we compared logistic regression on the raw 64-point window with logistic regression on the shift-invariant Fourier descriptors, on the windows as stored and after a random circular shift of every window (20 repeats). Robustness of the 1D-CNN, sparse logistic regression and the \textit{v}/\textit{c} rule was evaluated by training on clean data and testing on held-out folds with additive Gaussian noise at 30--0\,dB SNR or after decimation, over 10 repeats. Reporting follows TRIPOD+AI \cite{collins2024tripod}.

\section*{Data availability}
\authnote{State availability of the de-identified sensor and ECG data (for example, a public repository with DOI, or on reasonable request subject to ethics approval).}

\section*{Code availability}
All analysis code (data curation, feature extraction, models, evaluation and figures) is provided in the \texttt{jvp\_ml} directory and will be deposited in a public GitHub repository under the MIT License upon publication \authnote{add repository URL/DOI}.

\section*{Acknowledgements}
The authors thank the Science and Engineering Research Board (SERB), Government of India, for funding this work (File ID: SRG/2023/000732). \authnote{confirm funding applies to this project}

\section*{Author contributions}
\authnote{e.g. A.P. fabricated the sensor and acquired the data; A.C. ...; S.D. supervised the work. All authors reviewed the manuscript.}

\section*{Competing interests}
\authnote{Declare competing interests or state ``The authors declare no competing interests.''}

\bibliographystyle{naturemag}
\bibliography{jvp_references}

%% ======================================================================
\clearpage
\section*{Supplementary Information}
\setcounter{table}{0}\renewcommand{\thetable}{S\arabic{table}}
\setcounter{figure}{0}\renewcommand{\thefigure}{S\arabic{figure}}

\begin{table}[h]
\centering\footnotesize
\caption{\textbf{Data-curation issues found in the source workbooks} and how they were handled.}
\label{tab:qc}
\begin{tabularx}{\textwidth}{lX X}
\toprule
Issue & Participants (files) & Handling \\
\midrule
{{QCTABLE}}
\bottomrule
\end{tabularx}
\end{table}

\begin{figure}[h]
\centering
\includegraphics[width=0.75\textwidth]{figS1_sensitivity.pdf}
\caption{\textbf{Sensitivity analyses.} Cross-validated AUROC of core models (including the primary sensor 1D-CNN and sparse logistic regression) in the primary analysis, after excluding the five participants with copied workbooks (HS31, HS42, HS43, HS48, HS52), when using raw current instead of calibrated force, and when using only the first cycle per participant.}
\end{figure}

\begin{figure}[h]
\centering
\includegraphics[width=\textwidth]{figS2_ecg_investigation.pdf}
\caption{\textbf{Why the ECG benchmark is not a like-for-like comparison.} \textbf{a}, An ECG file (HS16, lead V1) and its six-Gaussian fit; the beats are parametric reconstructions. \textbf{b}, PR intervals measured on the reconstructed V1 beats; grey band, normal range (120--200\,ms); dashed line, first-degree AV-block threshold. \textbf{c}, A crude, start-agnostic \textit{v}/\textit{c} ratio computed from the raw sensor samples (windows rotated to their deepest sample; no interpolation or fiducial detection) against the physicians' diagnosis; dashed line, threshold of 1 ({{EI:agree}} of {{N}} participants classified concordantly).}
\end{figure}

\begin{table}[h]
\centering\scriptsize
\caption{\textbf{Class-imbalance correction strategies} (applied inside training folds). Sensitivity (Sens.) and specificity (Spec.) at a threshold of 0.5 and at the in-fold Youden threshold; CITL, calibration-in-the-large (mean predicted risk minus prevalence). LR, 20 repeats; RF, 5; CNN, 6 (threshold 0.5 only).}
\resizebox{\textwidth}{!}{%
\begin{tabular}{lllcccccccc}
\toprule
Input & Model & Strategy & AUROC & AUPRC & Sens.@0.5 & Spec.@0.5 & Sens.@in-fold & Spec.@in-fold & MCC@in-fold & CITL \\
\midrule
{{IMBTABLE}}
\bottomrule
\end{tabular}}
\end{table}

\begin{figure}[h]
\centering
\includegraphics[width=\textwidth]{figS3_imbalance.pdf}
\caption{\textbf{Effect of class imbalance (prevalence {{PREV}}).} \textbf{a}, AUROC and AUPRC of logistic regression on sensor (violet) and ECG (green) features under four strategies; dotted line, chance AUPRC. \textbf{b}, Sensitivity at a threshold of 0.5 (dotted) and at a threshold chosen inside the training folds (solid). \textbf{c}, Calibration-in-the-large. \textbf{d}, Sensor 1D-CNN under four training strategies. \textbf{e}, PPV projected over prevalence from the cross-validated sensitivity and specificity.}
\end{figure}

\end{document}
